% !TeX root = ../main.tex
\section{Feature Engineering and Selection}
\label{sec:features}

These features track battle trends and resource advantages that may influence
the final outcome.

\begin{itemize}[leftmargin=*,itemsep=0pt,parsep=0pt,topsep=2pt]
    \item \textbf{Base stats:} summarize team combat strength.
    \item \textbf{Team type weakness:} expose vulnerable matchups.
    \item \textbf{HP loss:} measure damage and momentum.
    \item \textbf{Type multipliers:} quantify attack effectiveness.
    \item \textbf{Statuses:} track damage, stat penalties, and lost turns.
    \item \textbf{Boosts:} measure changes in combat strength.
    \item \textbf{Effects:} capture temporary battle advantages.
    \item \textbf{Last-turn snapshot:} summarize the current state.
    \item \textbf{Faints:} P1 actual faints and P2 estimated faints.
    \item \textbf{Held items:} capture healing and damage bonuses.
    \item \textbf{Deployment:} counts, levels~\cite{pokemon_stat_formula}, effective HP reserves.
    \item \textbf{Tactics:} tempo/power, hazards, phazing, RestTalk~\cite{bulbapedia_metagame}.
\end{itemize}

\noindent\textbf{Feature selection.}
We first ran all features and retained those most correlated with the
win/loss target. For the three submitted logistic-regression/ensemble
configurations, we searched subsets by iterative SHAP-guided
pruning, repeatedly discarding weak predictors and reevaluating retained
features to obtain the final sets.